%%%%%%%%%%%%%%%%%%%%%%% file template.tex %%%%%%%%%%%%%%%%%%%%%%%%%
%
% This is a general template file for the LaTeX package SVJour3
% for Springer journals.          Springer Heidelberg 2010/09/16
%
% Copy it to a new file with a new name and use it as the basis
% for your article. Delete % signs as needed.
%
% This template includes a few options for different layouts and
% content for various journals. Please consult a previous issue of
% your journal as needed.
%
%%%%%%%%%%%%%%%%%%%%%%%%%%%%%%%%%%%%%%%%%%%%%%%%%%%%%%%%%%%%%%%%%%%
%
% First comes an example EPS file -- just ignore it and
% proceed on the \documentclass line
% your LaTeX will extract the file if required
\begin{filecontents*}{example.eps}
%!PS-Adobe-3.0 EPSF-3.0
%%BoundingBox: 19 19 221 221
%%CreationDate: Mon Sep 29 1997
%%Creator: programmed by hand (JK)
%%EndComments
gsave
newpath
  20 20 moveto
  20 220 lineto
  220 220 lineto
  220 20 lineto
closepath
2 setlinewidth
gsave
  .4 setgray fill
grestore
stroke
grestore
\end{filecontents*}
%
\RequirePackage{fix-cm}
%
%\documentclass{svjour3}                     % onecolumn (standard format)
%\documentclass[smallcondensed]{svjour3}     % onecolumn (ditto)
\documentclass[smallextended]{svjour3}       % onecolumn (second format)
%\documentclass[twocolumn]{svjour3}          % twocolumn
%
\smartqed  % flush right qed marks, e.g. at end of proof
%
\usepackage{graphicx}
\usepackage{cite}
\usepackage{amsmath}
\usepackage{mathtools}
\usepackage{multirow}
%\usepackage{algpseudocode}
\usepackage{algorithm}
\usepackage{hyperref}
\usepackage[noend]{algorithmic}
%\usepackage[]{algorithm2e}
\usepackage[normalem]{ulem}
 \useunder{\uline}{\ul}{}


\newcommand{\DE}{{\sc de}}
\newcommand{\DEEDM}{{\sc de-edm}}
\newcommand{\EAS}{{\sc ea}s}
\newcommand{\EA}{{\sc ea}}
\newcommand{\MULTIDYN}{{\sc multi-dyn}}
%
%%
%\usepackage{natbib}
%
% \usepackage{mathptmx}      % use Times fonts if available on your TeX system
%
% insert here the call for the packages your document requires
%\usepackage{latexsym}
% etc.
%
% please place your own definitions here and don't use \def but
% \newcommand{}{}
%
% Insert the name of "your journal" with
% \journalname{myjournal}
%
\begin{document}

\title{Differential Evolution with Enhanced Diversity Maintenance%\thanks{Grants or other notes
%about the article that should go on the front page should be
%placed here. General acknowledgments should be placed at the end of the article.}
}
%\subtitle{Do you have a subtitle?\\ If so, write it here}

%\titlerunning{Short form of title}        % if too long for running head

\author{Joel Chac\'on Castillo        \and
        Carlos Segura %etc.
}

%\authorrunning{Short form of author list} % if too long for running head

\institute{Joel Chac\'on Castillo \at
							Centro de Investigaci\'{o}n en Matemáticas, Guanajuato, Mexico\\
              \email{joel.chacon@cimat.mx}           %  \\
%             \emph{Present address:} of F. Author  %  if needed
           \and
					Carlos Segura \at
					Centro de Investigaci\'{o}n en Matemáticas, Guanajuato, Mexico\\
					\email{carlos.segura@cimat.mx}
}

\date{Received: date / Accepted: date}
% The correct dates will be entered by the editor


\maketitle

\begin{abstract}
Differential Evolution (\DE{}) is a popular population-based metaheuristic that has been successfully used in complex 
optimization problems.
%
Premature convergence is one of the most important drawbacks that affects its performance.
%
In combinatorial optimization, an alternative to alleviate premature convergence in memetic algorithms has been to 
explicitly control the diversity of the population.
%
However, this choice it is not usually applied in continuous optimization.
%
In this paper, a recently proposed replacement strategy to preserve diversity is extended and integrated with \DE{}
to generate the \DE{} with Enhanced Diversity Maintenance (\DEEDM{}).
%
The main novelty is the use of a dynamic balance between exploration and exploitation to adapt the proposal to the 
requirements of the different optimization stages.
%
Experimental validation is carried out with several benchmark tests proposed in competitions of the well-known IEEE Congress on 
Evolutionary Computation.
%
Top-rank algorithms of each competition are used to illustrate the usefulness of the proposal.
%
The new method avoids premature convergence and significantly improves further the results obtained by state-of-the-art algorithms.

\keywords{Diversity \and Differential Evolution \and Premature Convergence}
% \PACS{PACS code1 \and PACS code2 \and more}
% \subclass{MSC code1 \and MSC code2 \and more}
\end{abstract}

\section{Introduction}
\label{sec:Introduction}
Evolutionary Algorithms (\EAS{}) are one of the most widely used techniques to deal with complex optimization problems.
%
Several variants of these strategies have been devised~\cite{Talbi:09} and applied in many fields, such as in science, 
economic and engineering~\cite{chakraborty2008advances}.
%
Among them, Differential Evolution (\DE{}) \cite{storn1997differential} is one of the most effective strategies to deal
with continuous optimization.
%
In fact, it has been the winning strategy of several optimization competitions~\cite{das2011differential}.
%
Similarly to other \EAS{}, \DE{} is inspired by the natural evolution process and it involves the application of mutation, recombination and 
selection.
%
The main peculiarity of \DE{} is that it considers the differences among vectors that are present in the population to explore the search space.
%
In this sense it is similar to the Nelder-Mead~\cite{nelder1965simplex} and the 
Controlled Random Search (CRS) \cite{price1983global} optimizers.

In spite of the effectiveness of \DE{}, there exists several weaknesses that have been detected and partially
solved by extending the standard variant~\cite{das2011differential}.
%
Among them, the sensitivity to its parameters~\cite{zhang2009jade}, the appearance of stagnation due to the reduced exploration 
capabilities~\cite{sa2008exploration,lampinen2000stagnation} and premature convergence~\cite{zaharie2003control} are some of the most well-known
issues.
%
This last one issue is tackled in this paper.
%
Note that, attending to the proper design of population-based meta-heuristics~\cite{Talbi:09}, special attention must be
paid to attain a proper balance between exploration and exploitation.
%
A too large exploration degree prevents the proper intensification of the best located regions, usually resulting in a
too slow convergence.
%
Differently, an excessive exploitation degree provokes loss of diversity meaning that only a limited number of regions are sampled.
%
In the case of \DE{}, since its inception some criticism appeared because of its incapability to maintain a large
enough diversity due to the use of a selection with high pressure~\cite{sa2008exploration}.
%
Thus, several extensions of \DE{} to deal with premature convergence have been devised 
such as parameter adaptation~\cite{zaharie2003control}, 
auto-enhanced population diversity~\cite{yang2015differential} and selection strategies 
with a lower selection pressure~\cite{sa2008exploration}.
%
Some of the last studies on design of population-based meta-heuristics~\cite{Crepinsek:13} show that 
explicitly controlling the diversity to properly balance the exploration and intensification degree
is particularly useful.
%
Specifically, in the field of combinatorial optimization some novel replacement strategies that dynamically alter 
the balance between exploration and exploitation 
have appeared~\cite{segura2016novel}.
%
The main principle of such proposals is to use the stopping criterion and elapsed generations to bias the decisions 
taken by the optimizers with the aim of promoting exploration in the initial stages and exploitation in the last ones.
%
Probably their main weakness is that the time required to obtain high-quality solution increases.
%
Our novel proposal, which is called \DE{} with Enhanced Diversity Maintenance (\DEEDM{}), integrates a similar principle into \DE{}.
%
However, in order to avoid the excessive growth of computational requirements typical of diversity-based replacement strategies, 
two modifications that induce a larger degree of intensification are included.

The rest of the paper is organized as follows.
%
Some basic concepts of \DE{} and a review of works related to diversity within \DE{} are given in section~\ref{sec:Literature}.
%
Section~\ref{sec:trends} presents an analysis about the algorithms with best performance on the last continuous optimization 
contests held at the \textsc{ieee} Congress on Evolutionary Computation.
%
More emphasis is given on the variants based on \DE{}.
%
Our proposal is described in section~\ref{sec:Proposal}.
%
The experimental validation, which includes comparisons against state-of-the-art approaches, is detailed in section~\ref{sec:Experimental}. 
%
Finally, our conclusions and some lines of future work are given in section~\ref{sec:Conclusion}.


\section{Literature Review}
\label{sec:Literature}
\subsection{Differential Evolution: Basic Concepts}

This section is devoted to summarize the classic \DE{} variant and to introduce some of the most important terms used in the \DE{} field.
%
The classic \DE{} scheme is called the \DE{}/rand/1/bin and it has been extensively used to generate more complex \DE{} variants~\cite{das2011differential}.
%
In fact, our proposal also extends the classic variant.
%
%However, in order to perform fair comparisons,
%our experimental validation takes into account state-of-the-art approaches that incorporate more complex 
%components and even algorithms not belonging to the \DE{} field.
%
\DE{} was originally proposed as a direct search method for single-objective continuous optimization.
%
The variables governing a given problem performance are given as a vector like $\vec{X} = [x_1, x_2, ..., x_D]$, where $D$ is the
dimension of the problem.
%
In continuous optimization, each $x_i$ is a real number and usually box-constraints are given, i.e. there is a lower bound ($a_{i}$) and
upper bound ($b_{i}$) for each variable.
%
The aim of the optimization process is to obtain the vector $\vec{X}^*$ which minimizes a given objective function, mathematically 
denoted by $f : \Omega \subseteq \Re^D \rightarrow \Re$.
%
In the box-constrained case $\Omega = {\prod}_{j=1}^{D} [a_{j}, b_{j}]$.

\DE{} is a population-based stochastic algorithm, so it iteratively evolves a multi-set of candidate solutions.
%
In \DE{} such candidate solutions are usually called vectors.
%
In the basic \DE{} variant for each member of the population --- they are called \textit{target vectors} --- a new \textit{mutant vector}
is created.
%
Then, the mutant vector is combined with the target vector to generate a \textit{trial vector}.
%
Finally, a selection phase is applied to choose the survivors.
%
In this way, several generations are evolved until a stopping criterion is reached.
%
The $i$-th vector of the population at the generation $G$ is denoted as $\vec{X}_{i,G} = [x_{1,i,G}, x_{2,i,G},..., x_{D,i, G}]$.
%
In the following more details are given for each component of \DE{}.


\subsubsection{Initialization}

\DE{} usually starts the optimization process with a randomly initiated population of $N$ vectors.
%
Since there is commonly no information about the performance of different regions, uniform random generators are usually applied.
%
Hence, the $j$-th component of the $i$-th vector is initialized as $x_{j,i,0} = a_{j} + rand_{i,j}[0,1] (b_{j} - a_{j})$,
where $rand_{i,j}[0,1]$ is an uniformly distributed random number lying between $0$ and $1$.

\subsubsection{Mutation}

For each target vector a mutant vector is created and several ways of performing
such a process have been proposed.
%
In the classic \DE{} variant the rand/1 strategy is applied.
%
In this case, the mutant vector $V_{i,G}$ is created as follows:

\begin{equation}\label{eqn:mutation}
\vec{V}_{i,G} = \vec{X}_{r_1, G} + F \times (\vec{X}_{r_2, G} - \vec{X}_{r_3, G}) \quad r_1 \neq r_2 \neq r_3
\end{equation}
%
The indices $r_1, r_2, r_3 \in [1,N]$ are mutually different integers randomly chosen from the range $[1, N]$.
%
In addition, they are all different from the index $i$.
%
It is important to take into account that the difference between vectors is scaled with the number F, which is usually defined in the interval $[0.4, 1]$.
%
The scaled difference is added to a third vector, meaning that
when diversity decreases and consequently differences are low, mutant vectors are similar to target vectors.
%
As a result, maintaining some degree of diversity is specially important in \DE{}.

\subsubsection{Crossover}

In order to combine information of different candidate solutions and with the aim of increasing diversity, the crossover
operator is applied.
%
Specifically, each target vector $\vec{X}_{i,G}$ is mixed with its corresponding mutant vector $V_{i,G}$ to 
generate the trial vector $\vec{U_{i,G}} = [u_{1,i,G},u_{2,i,G}, ..., u_{D,i,G} ]$.
%
The most typical crossover is the \textit{binomial} one, which operates as follows:
%
\begin{equation} \label{eqn:crossover}
\vec{U}_{j,i,G}= 
\begin{cases}
    v_{j,i,G},& \text{if} (rand_{i,j}[0,1] \leq CR \quad or \quad j = j_{rand}  )\\
    v_{j,i,G},              & \text{otherwise}
\end{cases}
\end{equation}
where $rand_{i,j}[0,1]$ is a uniformly distributed random number,
$j_{rand}$ is a randomly chosen index which ensures that $\vec{U}_{i,G}$ inherits at least one component from $\vec{V}_{i,G}$ and
$CR \in [0,1]$ is the crossover rate.


\subsubsection{Selection}
Finally, a greedy selection is performed to determine the survivors of the next generation.
%
Each trial vector is compared with its corresponding target vector and the best one survives:

\begin{equation} \label{eqn:selection}
\vec{X}_{i,G+1}= 
\begin{cases}
    \vec{U}_{i,G},& \text{if} \quad f(\vec{U}_{i,G}) \leq f(\vec{X}_{i,G})  \\
    \vec{X}_{i,G},              & \text{otherwise}
\end{cases}
\end{equation}

Hence, each population member either gets better or remains with the same objective value in each generation.
%
Since members never deteriorate, it is considered to be a selection with high pressure.
%
Note that in case of a tie, the trial vector survives.

%
%The general convention is DE/\textit{x}/\textit{y}/\textit{x}, where DE indicates ``Differential Evolution'', \textit{x} denotes the base vector to be perturbed, \textit{y} is the number of difference vectors considered for perturbation and \textit{z} is the type of crossover to use.



\subsection{Diversity Preservation in Evolutionary Algorithms}

Given that premature convergence is a common drawback in  \EAS{}, several variants to deal with this issue have been designed \cite{eshelman1993real}.
%
Some of them start with a set of diverse candidate solutions, and as the generations evolve, such a diversity is reduced \cite{Crepinsek:13}.
%
Depending on the component of the \EA{} that is modified, these methods are usually classified in one the following groups~\cite{Crepinsek:13}: \textit{selection-based}, \textit{population-based}, \textit{crossover/mutation-based}, \textit{fitness-based}, and \textit{replacement-based}.
%
Recently, the replacement-based strategies have attained quite good performance, therefore this section is devoted to this kind
of methods.
%
Particularly, two different strategies that are used to validate our proposal are discussed.
%
One of the first techniques categorized as replacement-based are the \textit{crowding} methods.
%
A quite popular popular realization is the \textit{Restricted Tournament Selection} \cite{harik1995finding} (\RTS{}) strategy.
%
In \RTS{} after a new individual ($C$) is created, then $CF$ individuals from the current population are randomly selected.
%
Then, $C$ and its most similar individual --from those in the selected set-- compete for a place in the population using a traditional binary tournament, i.e. the best one survives.
%

Other strategies are based on considering the diversity to calculate a fitness value for the replacement stage.
%
Specifically, in the \textit{Hybrid Genetic Search with Adaptive Diversity Control} (\HGSADC{}) \cite{vidal2013hybrid}, the individuals --union of parents and offspring-- are sorted by their contribution to diversity and by their original cost.
%
In order to calculate the contribution to diversity of an individual, its mean distance to the closest 
$N_{Close}$ individuals is calculated.
%
For an individual $I$, the ranking in terms of diversity is denoted as $RD(I)$, whereas the ranking for the original cost is denoted as $RC(I)$.
%
Then, the rankings are combined to generate the biased fitness value $BF(I)$ using the Eq. \ref{eqn:biased}.
%
In each step of the replacement phase, the individual with the lowest biased fitness is erased and the ranks are re-calculated.
%
This process is performed until the desired population size is attained.
%
It is important to remark that this scheme requires the setting of the parameters: $N_{Close}$ and $N_{Elite}$,
whereas $N_{population}$ refers to the number of individuals that has not been erased yet by the replacement scheme.

\begin{equation}\label{eqn:biased}
 BF(I) = RC(I) + \left(1- \frac{N_{Elite}}{N_{population}} \right) RD(I)
\end{equation}
\subsection{Diversity in Differential Evolution}

\DE{} is highly susceptible to the loss of diversity due to the greedy strategy applied in the selection phase.
%
Thus, several analyses to better deal with this issue have been carried out.
%
Since the general implications of each \DE{} parameter on the diversity are known, one of
the alternatives is to theoretically estimate proper values for the \DE{} parameters~\cite{zaharie2003control}.
%
Differently, some analyses regarding the effects of the magnitude of the difference vectors used in the mutation
%Differently, some analyses regarding the effects of the Euclidean norm of the difference vectors used in the mutation
have also been performed~\cite{montgomery2009differential}.
%
Such analyses and additional empirical studies regarding the crossover allowed to conclude that some kind of movements 
should be disallowed to delay the convergence~\cite{montgomery2012simple}.
%
In this last study the kind of accepted movements varies along the run.
%
Specifically, it discards movements with a norm below a threshold and this threshold decreases taking into account the elapsed generations.
%
Other ways of altering the kind of accepted movements have been proposed~\cite{bolufe2013differential}.
%
Note that these kinds of methods have similarities with our proposal in the sense that decisions are biased by the number of elapsed generations.
%
However, our method operates on the replacement strategy and not on the mutation phase.
%
Moreover, these methods do not consider explicitly the differences appearing on the whole population.
%
Instead, the restrictions apply to the differences appearing in the reproduction phase.

A different alternative operates by altering the selection operator~\cite{sa2008exploration}.
%
Particularly, the selection pressure is relaxed through a probabilistic selection to maintain the population diversity and consequently 
to allow escaping from basin of attraction of local optima.
%
Since it considers the fitness to establish the acceptance probabilities, it is very sensitive to scale transformations.
%
In this case, decisions are not biased by the elapsed generations.

Finally, in the \textit{Auto-Enhanced Population Diversity} (\textsc{aepd}), the diversity is explicitly measured and it triggers a mechanism
to diversify the population when a too low diversity is detected~\cite{yang2015differential}.
%
Strategies with similar principles but with different disturbance schemes have also been devised~\cite{zhao2016differential}.

Note that \DE{} variants with best performance in competitions do not apply these modifications
and that most of these extensions have not been implemented in the most widely used optimization frameworks.
%
As a result, these extensions are not so widely used in the community in spite of their important benefits
for some cases.



\section{Performance in IEEE CEC Contests}
\label{sec:trends}
In the last decade, DE has been recognized as one of the most eficient and simple algorithm.
%
Particularly, the algorithms derivated from DE have obtained one of the top rank in the Congress on Evolutionary Computation (CEC).
%
In this work we are interested in the bias of DE algorithms and CEC competitions problems througth the last years.
%
It is important to take into account that althougth the kind of problems have changed, the DE are in the top rank.
%
However, in last years exist a tendency of the EAs, most of them considers the criteria stop as part of the adaptive mechanism (e.g. decreasing population size), and are derived as a mixture of several EA approaches (e.g. DE and Covariance Matrix Adaptation).
%

Some of the CEC competitions are described as follows:



\section{Proposal}
\label{sec:Proposal}
Our proposal is motivated by two main works in the area of diversity-based EAs.
%
First, the empirical studies developed by Montgomery et al~\cite{montgomery2012simple},
which presents several empirical analyses to confirm issues related to premature convergence in \DE{}.
%
The second work, by Segura et al.~\cite{segura2016novel}, provides significant improvements in the combinatorial optimization field.
%
Particularly, they incorporated a novel replacement strategy called \textit{The Replacement with Multi-objective based Dynamic Diversity Control} (\RMDDC{}) 
to relate the control of diversity with the stopping criterion and elapsed generations.
%
Important benefits were attained by methods including \RMDDC{}, so given the conclusions of these previous works, the proposal of this paper is a 
novel \DE{} variant that includes an explicit mechanism that follows similar principles to \RMDDC{}.
%
This novel optimizer is called \textit{Differential Evolution with Enhanced Diversity Maintenance} (\DEEDM{}) and its source
code is freely available~\footnote{The code in C++ can be downloaded in the next link \url{https://github.com/joelchaconcastillo/Diversity\_DE\_Research.git}}.

The novelty of \DEEDM{} (see Algorithm~\ref{alg:DEEDM}) is the incorporation of a replacement strategy that follows some of the principles of the one used in \RMDDC{}, i.e.
individuals that contribute too little to diversity should not be accepted as members of the next generation.
%
In this way, the greedy selection strategy of \DE{} is not used to maintain the parent population ($X$).
%
However, an additional population, which is called the elite population ($E$), is incorporated and this one is managed in a similar way than 
the population of \DE{}, i.e. with a greedy replacement.
%
Moreover, in order to establish the minimum acceptable diversity contribution to be selected, the stopping criterion and elapsed
generations are taken into account.
%
One of the main weaknesses of this method is that its convergence is highly delayed.
%
Thus, in order to promote a faster convergence than in \RMDDC{} two modifications are performed.
%
First, no concepts of the multi-objective field are applied, instead a more greedy selection is taken into account.
%
Second, the elite population is also considered as an input of the replacement strategy.

\begin{algorithm}[t]
\algsetup{linenosize=\tiny}
  \scriptsize
	\caption{General scheme of DE-EDM} 
	\begin{algorithmic}[1]
	\STATE Randomly initialize the population of $NP$ individuals, where each one is uniformly distributed.
	\STATE $G=0$
	\WHILE{ stopping criterion is not satisfied}
	   \FOR{ $i=1$ to $NP$} 
		\STATE Mutation: Generate the donor vector ($V_{i,G}$) according Eq. (\ref{eqn:mutation}).
		\STATE Crossover: Recombine the mutate vector ($U_{i,G}$) according Eq. (\ref{eqn:crossover}).
		\STATE Selection: Update the elite vector ($E_{i,G}$ instead of $X_{i,G}$) according Eq. (\ref{eqn:selection}).
		\STATE Replacement: Select the target vectors ($X_{i,G+1}$) according to Algorithm \ref{alg:Replacement} .
	   \ENDFOR
	   \STATE $G=G+1$
	\ENDWHILE
    \label{alg:DEEDM}
\end{algorithmic}
\end{algorithm}

%\DEEDM{} alters the replacement strategy of \DE{} with the aim of controlling the 
%balance between exploration and exploitation by extending \RMDDC{}.
%
%The principles of the \RMDDC{} are as follows.
%
%It considers the maximization of the diversity contribution of each individual as an explicit objective.
%
%Then, the notion of Pareto dominance is used to select the survivors.
%
%It uses a dynamic threshold to prevent the selection of individuals with low contribution to diversity.
%
%
%Thus, executions of several days were required to attain high-quality solutions in many of the tested problems.
%
%As a result our proposal incorporates two extensions to alleviate such a drawback.
%
%This modification considers the inclusion of an elite population, which records the individuals with the best fitness.
%

Our replacement strategy (see Algorithm \ref{alg:Replacement}) operates as follows.
%
It receives as input the parent population (target vectors), the offspring population (trial vectors), and the elite population.
%
In each generation it must select the $NP$ vectors of the next parent population.
%
First, it calculates the desired minimum distance $D_t$ given the current number of elapsed function evaluations (line 2).
%
Then, it joins the three populations in a set of current members (line 3).
%
The current members set contains vectors that might be selected to survive.
%
Then, the set of survivors and penalized individuals are initialized to the empty set (line 4).
%
In order to select the $NP$ survivors (next parent population) an iterative process is repeated (lines 5 - 13).
%
In each step the best individual in the \textit{Current set}, i.e. the one with best objective function is selected
to survive, i.e. it is moved to the \textit{Survivor set} (line 6 - 8).
%
Then, individuals in the \textit{Current set} with a distance metric lower than $D_t$ are transferred to the \textit{Penalized set} (line 9).
%
%%The way to calculate the distance between two individuals is by a using a normalized Euclidean distance (Eq.~\ref{eqn:distance}).
The way to calculate the distance between two individuals is by using the normalized Euclidean distance decribed in the equation~\ref{eqn:distance}, where $D$ is the dimension of the problem, and $a_d, b_d$ are the minimum and maximum bounds of each dimension ($d$).
%
%
In cases where the \textit{Current set} is empty previous to the selection of $NP$ individuals, the \textit{Survivor set} is filled by selecting in each step 
the individual in $Penalized$ with the largest distance to the closest individual in the \textit{Survivor set}.

\begin{equation}\label{eqn:distance}
distance ( x_{i}, x_j ) = \frac{\sqrt{ \sum_{d=1}^D \left ( \frac{x_{i}^d - x_j^d}{b_d - a_d} \right )^2  }} {\sqrt{D}}
\end{equation}


\begin{algorithm}[t]
\algsetup{linenosize=\tiny}
  \scriptsize
	\caption{Replacement Phase} \label{alg:Replacement}
	\begin{algorithmic}[1]
	\STATE Input: $Population$ ($target$ $vectors$), $Offspring$ ($trial$ $vectors$), and $Elite$
	\STATE Update $D_t = D_I - D_I *(nfes/(0.95*max\_nfes)) $ 
	\STATE $Current = Population \cup Offspring \cup Elite$.
	\STATE $Survivors = Penalized = \emptyset$.
	\WHILE{ $|Survivors| < NP$ And $|Current| > 0$ }
	   \STATE $Selected$ = Select the best individual of $Current$.
		 \STATE Remove $Selected$ from $Current$.
	   \STATE Copy $Selected$ to $Survivors$.
	   \STATE Find the individuals from $Current$ with a distance to $Selected$ lower than $D_t$ and move them to $Penalized$. Normalized distance is considered (Eq. \ref{eqn:distance}).
	\ENDWHILE
	\WHILE{ $|Survivors| < NP$ }
	   \STATE $Selected$ = Select the individual from $Penalized$ with the largest distance to the closest individual in $Survivors$.
		 \STATE Remove $Selected$ from $Penalized$.
	   \STATE Copy $Selected$ to $Survivors$.
	\ENDWHILE
  \RETURN $Survivors$
\end{algorithmic}
\end{algorithm}


In order to complete the description is important to specify the way to calculate $D_t$ and the methodology to update the 
elite individuals.
%
All the remaining steps are maintained as in the classic \DE{} variant.
%
The value of $D_t$ is used to alter the degree between exploration and explotation so it should depend on the optimization stage.
%
Specifically, this value should be reduced as the stopping criterion is reached with the aim of promoting explotation.
%
In our scheme, an initial value for $D_t$ ($D_I$) must be set.
%
Then, similarly than in~\cite{segura2016novel}, a linear reduction of $D_t$ is performed by taking into account the elepased function evaluations and stopping criterion.
%
Particularly, in this work, the stopping criterion is set by function evaluations (\textit{nfes}).
%
The reduction is calculated in such way that by the $95\%$ of maximum number of evaluations the resulting $D_t$ value is $0$.
%
Therefore, in the remaining $5\%$ diversity is not considered at all.
%
Thus, if $max\_nfes$ is the maximum number of evaluations and $nfes$ is the elapsed number of evaluations, $D_t$ can be calculated as $D_t=D_I - D_I *(nfes/(0.95*max\_nfes))$.
%
%The reduction of this parameter is based in a linear reduction, as is indicated in .
%
%They indicated that a linear decrement generally provides the most stable results.
%%
%%We consider a linear model reduction based in empirical studies of similar works \cite{segura2016novel}, where they indicated that a linear decrement generally provides the most stable results.
%%%

%Principally, the replacement phase considers a defined radius ($D_t$), which for simplicity is considered through the normalized euclidean distance (Eq. \ref{eqn:distance}), however other distances could be implemented (e.g. mahalanobis distance).
%
%It is important to mention that the normalized distance is needed, thus each dimension is equally important, also this simplifies the setting of the $D_I$ parameter, since that the maximum distance is the unity, which is a fraction of the main space diagonal.


The initial distance ($D_I$) heavily affects the performance of \DEEDM{}.
%
If this parameter is fixed large enough, then at the first optimization stages the algorithm aims to maximize the diversity of the population, 
so a proper exploration is performed which is very important in some kinds of problems such as highly multimodal and deceptive ones.
%
Thus, the effect of premature convergence might be alleviated.
%
A too large $D_I$ might induce too much exploration so a proper exploitation phase is not performed.
%
In the opposite case, a too low $D_I$ might avoid the exploration phase, so avoiding local optima is more difficult.
%
Depending on the kind of fitness landscape and stopping criterion, the optimal $D_I$ might vary.
%
For instance, deceptive and highly multi-modal problems usually require larger values than unimodal problems.
%
However, in our proposal, $D_I$ is not adapted to each problem, instead an experimental study to check the robustness
of different $D_I$ value is attached in the experimental validation section. 
%
%It is interesting to take into account that our proposal could be translated to the multi-modal fashion just setting a final distance factor, which should be strictly positive, however this trend is not considered in this work.
%
%Therefore, when the initial distance factor is setted several factors should be taken into account.
%
%One of them is the maximum number of evaluations, when the problem is setted with few number of evaluations, the initial distance factor should be low, since the exploitation should be promoted.
%
%On the other hand, when the problem is configured with a large number of evaluations, then the initial distance factor should be higher.
%
%Also, the number of vectors (population size) should be considered, since that they are directly related with the diversity.
%
%Although our proposal should take into account the previously details, there exist an usual range where the results are enough stable, it is showed in the empirical analyzes.
%
%Our proposal provides several advantages in contrast of the standard DE and some of the state-of-the-art algorithms, perhaps the most important is that it is not over-parameterized.
%
%Although several top-rank algorithms present good enough results, usually these algorithms need a tuning phase to fit the parameters, therefore in real application it could be infeasible.

%
Similarly that the standard \DE{}, in \DEEDM{} the crossover probability ($CR$) and the mutation factor ($F$) must be set.
%
The first one is perhaps the most important for the performance according to several studies developed by Montgomery 
et al. \cite{montgomery2010analysis}.
%
These authors empirically proved that extremes $CR$ values leads to vastly different search behaviors.
%
They explained that low $CR$ values result in a search that is aligned with a small number of search space axes and
induce small displacements.
%
This provokes a gradual and slow convergence that in some scenarios might result in a robust behavior.
%
Additionally, high $CR$ values might generate quality solutions with a lower probability.
%
However, these transformations provoke large displacements that could improve significantly the solutions.
%
According to this, we employ both high and low $CR$ values as it is showed in eq. \ref{eqn:cr}.

\begin{equation} \label{eqn:cr}
CR = 
\begin{cases}
     Normal(0.2, 0.1),& \text{if} \quad rand[0,1] \leq 0.5  \\
     Normal(0.9, 0.1),              & \text{otherwise}
\end{cases}
\end{equation}


%TODO: reescribir
Following the principles of several SHADE variants~\cite{awad2016ensemble, brest2016shade}, the function evaluations are considered in the random generation of the mutation factor $F$.
%
Particularly, each $F$ is sampled through a Cauchy distribution (eq. \ref{eqn:cauchy}).
\begin{equation}\label{eqn:cauchy}
 Cauchy(0.5, 0.5*nfes/max\_nfes)
\end{equation}
%
Therefore, at the first optimization stages are generated values near of $0.5$.
%
However, as the execution transcurs, the density function suffer a gradual transformation, thus the variance is increased.
%
It shows a bias to sample values outside the interval $[0.0, 1.0]$.
%
Where the values larger than $1.0$ are trucated, and the negative values are sampled again (until it generates a positive factor).
%
Particularly, the aim of this is to generate small $F$-values at the first optimization stages, and at the end larger $F$-values ($F=1$), in such way that the fast convergence can be avoided.
%
%

% thus the shape of the distribution increases with the function evaluations and therefore are generated more extreme values at the end of execution, this aims avoid stagnation in different stages of the algorithm.
%


%Thus, briefly our proposal is based in several ideas of the mentioned works, and are listed as follows:
%\begin{itemize}
%\item Is considered a threshold to control explicitly the convergence of the solutions.
%\item This threshold decreases over the algorithm's run.
%\item The selection operator is relaxed in the sense that it does not provokes premature convergence, this is attained considering an elite population. 
%\end{itemize}



\section{Experimental Study}
\label{sec:Experimental}
In order to statistically compare the hypervolume results, a similar guideline than the one proposed in~\cite{Joel:StatisticalTest} was used. 
%
First a Shapiro-Wilk test was performed to check whatever or not the values of the results followed a Gaussian distribution. 
%
If, so, the Levene test was used to check for the homogeneity of the variances. 
%
If samples had equal variance, an ANOVA test was done; if not, a Welch test was performed. 
%
For non-Gaussian distributions, the nonparametric Kruskal-Wallis test was used to test whether samples are drawn from the same distribution. 
%
An algorithm $X$ is said to win algorithm $Y$ when the differences between them are statistically significant, if the mean and median obtained by $X$ are higher than the mean and median achieved by $Y$.
% Please add the following required packages to your document preamble:
% \usepackage{multirow}
% \usepackage{graphicx}
\begin{table}[]
\centering
\caption{My caption}
\label{my-label}
\resizebox{\textwidth}{!}{%
\begin{tabular}{|c|c|c|c|c|c|c|}
\hline
\multirow{2}{*}{\textbf{Algorithm}} & \multirow{2}{*}{\textbf{Always Solved}} & \multirow{2}{*}{\textbf{\begin{tabular}[c]{@{}c@{}}At least one \\ time solved\end{tabular}}} & \multirow{2}{*}{\textbf{Score}} & \multicolumn{3}{c|}{\textbf{Statistical tests}} \\ \cline{5-7} 
 &  &  &  & $\uparrow$ & $\downarrow$ & $\longleftrightarrow$ \\ \hline
\textbf{EBOwithCMAR} & 11 & 15 & 30.6792 & 22 & 32 & 6 \\ \hline
\textbf{JSO} & 8 & 19 & 41.8322 & 25 & 29 & 6 \\ \hline
\textbf{Proposal} & 21 & 28 & 100.0000 & 36 & 22 & 2 \\ \hline
\end{tabular}%
}
\end{table}

% Please add the following required packages to your document preamble:
% \usepackage{multirow}
% \usepackage{graphicx}
\begin{table}[]
\centering
\caption{My caption}
\label{my-label}
\resizebox{\textwidth}{!}{%
\begin{tabular}{|c|c|c|c|c|c|c|}
\hline
\multirow{2}{*}{\textbf{Algorithm}} & \multirow{2}{*}{\textbf{Always Solved}} & \multirow{2}{*}{\textbf{\begin{tabular}[c]{@{}c@{}}At least one \\ time solved\end{tabular}}} & \multirow{2}{*}{\textbf{Score}} & \multicolumn{3}{c|}{\textbf{Statistical tests}} \\ \cline{5-7} 
 &  &  &  & $\uparrow$ & $\downarrow$ & $\longleftrightarrow$ \\ \hline
\textbf{UMOEAsII} & 9 & 14 & 41.65 & 19 & 31 & 10 \\ \hline
\textbf{L-SHADE-Epsilon} & 7 & 13 & 45.84 & 25 & 23 & 12 \\ \hline
\textbf{Proposal} & 13 & 21 & 100.00 & 31 & 21 & 8 \\ \hline
\end{tabular}%
}
\end{table}

% Please add the following required packages to your document preamble:
% \usepackage{graphicx}
\begin{table}[]
\centering
\caption{My caption}
\label{my-label}
\resizebox{\textwidth}{!}{%
\begin{tabular}{|c|c|c|c|c|c|c|}
\hline
 & \textbf{Best} & \textbf{Worst} & \textbf{Median} & \textbf{Mean} & \textbf{Std} & \textbf{Succ. Ratio} \\ \hline
$f_1$ & 0.00E+00 & 0.00E+00 & 0.00E+00 & 0.00E+00 & 0.00E+00 & 1.00E+00 \\ \hline
$f_2$ & 0.00E+00 & 0.00E+00 & 0.00E+00 & 0.00E+00 & 0.00E+00 & 1.00E+00 \\ \hline
$f_3$ & 0.00E+00 & 0.00E+00 & 0.00E+00 & 0.00E+00 & 0.00E+00 & 1.00E+00 \\ \hline
$f_4$ & 0.00E+00 & 0.00E+00 & 0.00E+00 & 0.00E+00 & 0.00E+00 & 1.00E+00 \\ \hline
$f_5$ & 0.00E+00 & 0.00E+00 & 0.00E+00 & 0.00E+00 & 0.00E+00 & 1.00E+00 \\ \hline
$f_6$ & 0.00E+00 & 0.00E+00 & 0.00E+00 & 0.00E+00 & 0.00E+00 & 1.00E+00 \\ \hline
$f_7$ & 0.00E+00 & 0.00E+00 & 0.00E+00 & 0.00E+00 & 0.00E+00 & 1.00E+00 \\ \hline
$f_8$ & 0.00E+00 & 0.00E+00 & 0.00E+00 & 0.00E+00 & 0.00E+00 & 1.00E+00 \\ \hline
$f_9$ & 0.00E+00 & 0.00E+00 & 0.00E+00 & 0.00E+00 & 0.00E+00 & 1.00E+00 \\ \hline
$f_{10}$ & 0.00E+00 & 1.20E-01 & 0.00E+00 & 1.65E-02 & 3.39E-02 & 7.45E-01 \\ \hline
$f_{11}$ & 0.00E+00 & 0.00E+00 & 0.00E+00 & 0.00E+00 & 0.00E+00 & 1.00E+00 \\ \hline
$f_{12}$ & 0.00E+00 & 2.20E-01 & 0.00E+00 & 6.37E-02 & 1.76E-01 & 6.67E-01 \\ \hline
$f_{13}$ & 0.00E+00 & 0.00E+00 & 0.00E+00 & 0.00E+00 & 0.00E+00 & 1.00E+00 \\ \hline
$f_{14}$ & 0.00E+00 & 0.00E+00 & 0.00E+00 & 0.00E+00 & 0.00E+00 & 1.00E+00 \\ \hline
$f_{15}$ & 0.00E+00 & 0.00E+00 & 0.00E+00 & 0.00E+00 & 0.00E+00 & 1.00E+00 \\ \hline
$f_{16}$ & 0.00E+00 & 2.10E-01 & 0.00E+00 & 2.47E-02 & 7.27E-02 & 8.82E-01 \\ \hline
$f_{17}$ & 0.00E+00 & 0.00E+00 & 0.00E+00 & 0.00E+00 & 0.00E+00 & 1.00E+00 \\ \hline
$f_{18}$ & 0.00E+00 & 1.00E-02 & 0.00E+00 & 1.96E-03 & 4.47E-03 & 8.04E-01 \\ \hline
$f_{19}$ & 0.00E+00 & 0.00E+00 & 0.00E+00 & 0.00E+00 & 0.00E+00 & 1.00E+00 \\ \hline
$f_{20}$ & 0.00E+00 & 0.00E+00 & 0.00E+00 & 0.00E+00 & 0.00E+00 & 1.00E+00 \\ \hline
$f_{21}$ & 0.00E+00 & 0.00E+00 & 0.00E+00 & 0.00E+00 & 0.00E+00 & 1.00E+00 \\ \hline
$f_{22}$ & 0.00E+00 & 0.00E+00 & 0.00E+00 & 0.00E+00 & 0.00E+00 & 1.00E+00 \\ \hline
$f_{23}$ & 0.00E+00 & 3.00E+02 & 0.00E+00 & 3.49E+01 & 1.03E+02 & 8.82E-01 \\ \hline
$f_{24}$ & 0.00E+00 & 0.00E+00 & 0.00E+00 & 0.00E+00 & 0.00E+00 & 1.00E+00 \\ \hline
$f_{25}$ & 0.00E+00 & 1.00E+02 & 0.00E+00 & 3.92E+00 & 2.00E+01 & 9.61E-01 \\ \hline
$f_{26}$ & 0.00E+00 & 0.00E+00 & 0.00E+00 & 0.00E+00 & 0.00E+00 & 1.00E+00 \\ \hline
$f_{27}$ & 0.00E+00 & 3.87E+02 & 3.87E+02 & 2.05E+02 & 2.68E+02 & 1.96E-02 \\ \hline
$f_{28}$ & 0.00E+00 & 0.00E+00 & 0.00E+00 & 0.00E+00 & 0.00E+00 & 1.00E+00 \\ \hline
$f_{29}$ & 1.45E+02 & 2.26E+02 & 2.18E+02 & 1.99E+02 & 4.21E+01 & 0.00E+00 \\ \hline
$f_{30}$ & 3.95E+02 & 3.95E+02 & 3.95E+02 & 3.95E+02 & 2.10E-01 & 0.00E+00 \\ \hline
\end{tabular}%
}
\end{table}

% Please add the following required packages to your document preamble:
% \usepackage{graphicx}
\begin{table}[]
\centering
\caption{My caption}
\label{my-label}
\resizebox{\textwidth}{!}{%
\begin{tabular}{|c|c|c|c|c|c|c|}
\hline
 & \textbf{Best} & \textbf{Worst} & \textbf{Median} & \textbf{Mean} & \textbf{Std} & \textbf{Succ. Ratio} \\ \hline
$f_1$ & 0.00E+00 & 0.00E+00 & 0.00E+00 & 0.00E+00 & 0.00E+00 & 1.00E+00 \\ \hline
$f_2$ & 0.00E+00 & 0.00E+00 & 0.00E+00 & 0.00E+00 & 0.00E+00 & 1.00E+00 \\ \hline
$f_3$ & 0.00E+00 & 0.00E+00 & 0.00E+00 & 0.00E+00 & 0.00E+00 & 1.00E+00 \\ \hline
$f_4$ & 0.00E+00 & 0.00E+00 & 0.00E+00 & 0.00E+00 & 0.00E+00 & 1.00E+00 \\ \hline
$f_5$ & 0.00E+00 & 0.00E+00 & 0.00E+00 & 0.00E+00 & 0.00E+00 & 1.00E+00 \\ \hline
$f_6$ & 0.00E+00 & 3.60E-02 & 4.00E-03 & 7.39E-03 & 1.15E-02 & 3.92E-01 \\ \hline
$f_7$ & 2.00E-02 & 1.02E-01 & 5.90E-02 & 5.77E-02 & 4.93E-02 & 0.00E+00 \\ \hline
$f_8$ & 0.00E+00 & 0.00E+00 & 0.00E+00 & 0.00E+00 & 0.00E+00 & 1.00E+00 \\ \hline
$f_9$ & 0.00E+00 & 0.00E+00 & 0.00E+00 & 0.00E+00 & 0.00E+00 & 1.00E+00 \\ \hline
$f_{10}$ & 0.00E+00 & 0.00E+00 & 0.00E+00 & 0.00E+00 & 0.00E+00 & 1.00E+00 \\ \hline
$f_{11}$ & 0.00E+00 & 6.00E-02 & 0.00E+00 & 5.88E-03 & 1.90E-02 & 9.02E-01 \\ \hline
$f_{12}$ & 0.00E+00 & 0.00E+00 & 0.00E+00 & 0.00E+00 & 0.00E+00 & 1.00E+00 \\ \hline
$f_{13}$ & 1.00E-02 & 8.00E-02 & 5.00E-02 & 4.67E-02 & 2.60E-02 & 0.00E+00 \\ \hline
$f_{14}$ & 1.00E-02 & 5.00E-02 & 3.00E-02 & 2.82E-02 & 2.13E-02 & 0.00E+00 \\ \hline
$f_{15}$ & 0.00E+00 & 4.70E-01 & 2.20E-01 & 1.99E-01 & 1.55E-01 & 1.96E-02 \\ \hline
$f_{16}$ & 4.00E-02 & 1.50E-01 & 8.00E-02 & 8.47E-02 & 4.96E-02 & 0.00E+00 \\ \hline
$f_{17}$ & 0.00E+00 & 0.00E+00 & 0.00E+00 & 0.00E+00 & 0.00E+00 & 1.00E+00 \\ \hline
$f_{18}$ & 0.00E+00 & 2.00E-02 & 1.00E-02 & 7.65E-03 & 6.32E-03 & 3.14E-01 \\ \hline
$f_{19}$ & 0.00E+00 & 0.00E+00 & 0.00E+00 & 0.00E+00 & 0.00E+00 & 1.00E+00 \\ \hline
$f_{20}$ & 0.00E+00 & 0.00E+00 & 0.00E+00 & 0.00E+00 & 0.00E+00 & 1.00E+00 \\ \hline
$f_{21}$ & 0.00E+00 & 0.00E+00 & 0.00E+00 & 0.00E+00 & 0.00E+00 & 1.00E+00 \\ \hline
$f_{22}$ & 0.00E+00 & 3.00E-02 & 0.00E+00 & 3.73E-03 & 2.76E-02 & 7.65E-01 \\ \hline
$f_{23}$ & 0.00E+00 & 1.00E+02 & 0.00E+00 & 2.55E+01 & 5.10E+01 & 7.45E-01 \\ \hline
$f_{24}$ & 0.00E+00 & 6.90E-01 & 0.00E+00 & 2.61E-02 & 1.33E-01 & 9.61E-01 \\ \hline
$f_{25}$ & 1.00E+02 & 1.00E+02 & 1.00E+02 & 1.00E+02 & 0.00E+00 & 0.00E+00 \\ \hline
$f_{26}$ & 8.00E-02 & 1.00E+02 & 5.29E+01 & 5.20E+01 & 3.19E+01 & 0.00E+00 \\ \hline
$f_{27}$ & 2.50E-01 & 9.10E-01 & 5.40E-01 & 5.60E-01 & 2.92E-01 & 0.00E+00 \\ \hline
$f_{28}$ & 0.00E+00 & 3.57E+02 & 3.43E+02 & 2.76E+02 & 1.60E+02 & 1.96E-01 \\ \hline
$f_{29}$ & 1.00E+02 & 1.00E+02 & 1.00E+02 & 1.00E+02 & 0.00E+00 & 0.00E+00 \\ \hline
$f_{30}$ & 1.84E+02 & 1.84E+02 & 1.84E+02 & 1.84E+02 & 3.25E-02 & 0.00E+00 \\ \hline
\end{tabular}%
}
\end{table}


\section{Conclusion}
\label{sec:Conclusion}
Premature convergence is one of the most important drawbacks of \DE{}.
%
This paper proposes a novel variant of \DE{}, the \textit{Differential Evolution with Enhanced Diversity Maintenance} (\DEEDM{}),
which incorporates a novel replacement phase that is based on a recently proposed strategy (\RMDDC{}) that was successfully
used with memetic algorithms in the combinatorial optimization field.
%
In order to attain a faster convergence than in \RMDDC{} two modifications are proposed for the integration with \DE{}.
%
First, no concepts of the multi-objective field are applied, meaning that a more greedy strategy is devised.
%
Second, an elite population is incorporated.
%
Experimental validation shows the important benefits obtained by \DEEDM{} in long-term executions both when compared with a classic \DE{} variant and with
the state-of-the-art proposals.
%
\DEEDM{} requires setting a new parameter, which is called the initial distance factor.
%
The obtained quality highly depends on the proper specification of this parameter.
%
However, a fixed value could be used for properly dealing with benchmarks of the CEC 2016 and CEC 2017 competitions,
meaning that a quite robust proposal is obtained.

Several extensions might be explored to further improve our proposal.
%
First, with the aim of facilitating the application of \DEEDM{}, adaptive and/or self-adaptive schemes might be applied for setting the initial distance
factor, as well as for other \DE{} parameters.
%
Second, with the aim of reducing the number of evaluations required to attain high-quality solutions, some of the strategies that are used in the
field of optimization with expensive functions might be integrated into our proposal.
%
Finally, the incorporation of a local search engine mainly in the last stages of the optimization might bring additional benefits.


%Text with citations \cite{RefB} and \cite{RefJ}.
%%\subsection{Subsection title}
%%\label{sec:2}
%%as required. Don't forget to give each section
%%and subsection a unique label (see Sect.~\ref{sec:1}).
%%\paragraph{Paragraph headings} Use paragraph headings as needed.
%%% For one-column wide figures use
%%\begin{figure}
%%% Use the relevant command to insert your figure file.
%%% For example, with the graphicx package use
%%  \includegraphics{example.eps}
%%% figure caption is below the figure
%%\caption{Please write your figure caption here}
%%\label{fig:1}       % Give a unique label
%%\end{figure}
%
% For two-column wide figures use
%\begin{figure*}
%% Use the relevant command to insert your figure file.
%% For example, with the graphicx package use
%  \includegraphics[width=0.75\textwidth]{example.eps}
%% figure caption is below the figure
%\caption{Please write your figure caption here}
%\label{fig:2}       % Give a unique label
%\end{figure*}
%
%% For tables use
%\begin{table}
%% table caption is above the table
%\caption{Please write your table caption here}
%\label{tab:1}       % Give a unique label
%% For LaTeX tables use
%\begin{tabular}{lll}
%\hline\noalign{\smallskip}
%first & second & third  \\
%\noalign{\smallskip}\hline\noalign{\smallskip}
%number & number & number \\
%number & number & number \\
%\noalign{\smallskip}\hline
%\end{tabular}
%\end{table}

%\begin{acknowledgements}
%If you'd like to thank anyone, place your comments here
%and remove the percent signs.
%\end{acknowledgements}

% BibTeX users please use one of
%\bibliographystyle{spbasic}      % basic style, author-year citations
%\bibliographystyle{spmpsci}      % mathematics and physical sciences
\bibliographystyle{spphys}       % APS-like style for physics
%\bibliography{}   % name your BibTeX data base
\bibliography{example}                % name your BibTeX data base

% Non-BibTeX users please use
%\begin{thebibliography}{}
%%
%% and use \bibitem to create references. Consult the Instructions
%% for authors for reference list style.
%%
%\bibitem{RefJ}
%% Format for Journal Reference
%Author, Article title, Journal, Volume, page numbers (year)
%% Format for books
%\bibitem{RefB}
%Author, Book title, page numbers. Publisher, place (year)
%% etc
%\end{thebibliography}

\end{document}
% end of file template.tex
